% Kompilacja: pdfLaTeX + BibTeX + pdfLaTeX x2
% Przed redakcją przeczytaj ../analysis/ANALIZA.md i ../analysis/DECYZJA.md.
\documentclass[12pt,a4paper]{article}

\usepackage[T1]{fontenc}
\usepackage[utf8]{inputenc}
\usepackage[polish]{babel}
\usepackage[a4paper,margin=2.5cm]{geometry}
\usepackage{amsmath,amssymb}
\usepackage{graphicx}
\usepackage{booktabs}
\usepackage{siunitx}
\usepackage{caption}
\usepackage{float}
\usepackage[hidelinks]{hyperref}

\sisetup{output-decimal-marker={,},group-separator={\,},detect-weight=true}
\graphicspath{{../figures/}}

\title{Tytuł ćwiczenia}
\author{Karina Nepravskaya}
\date{Kraków \the\year}

\begin{document}
\maketitle

\section{Wstęp teoretyczny}
% Tylko teoria potrzebna do metody i interpretacji.

\section{Cel ćwiczenia}

\section{Aparatura}

\section{Przebieg doświadczenia}
% Opisz faktyczny przebieg w formie bezosobowej.

\section{Opracowanie wyników}
% Zdefiniuj symbole, podaj równania, pokaż obliczenia i wykresy.

\section{Niepewności pomiarowe}

\section{Wyniki}

\section{Wnioski i podsumowanie}

\begin{thebibliography}{9}
\bibitem{instrukcja}
Instrukcja do ćwiczenia laboratoryjnego, dostęp: dzień miesiąc rok.
\end{thebibliography}

\end{document}
